\documentclass[11pt]{article}
\usepackage[T1]{fontenc}
\usepackage[utf8]{inputenc}
\usepackage[margin=1in]{geometry}
\usepackage{amsmath,amssymb,amsthm}
\usepackage{enumitem}
\usepackage{booktabs,array}
\usepackage[colorlinks=true,linkcolor=blue,citecolor=blue,urlcolor=blue]{hyperref}
\usepackage{microtype}
\newtheorem{theorem}{Theorem}[section]
\newtheorem{lemma}[theorem]{Lemma}
\newtheorem{proposition}[theorem]{Proposition}
\newtheorem{corollary}[theorem]{Corollary}
\theoremstyle{definition}
\newtheorem{definition}[theorem]{Definition}
\theoremstyle{remark}
\newtheorem{remark}[theorem]{Remark}
\newcommand{\R}{\mathbb R}
\newcommand{\AZ}{\mathrm{AZ}}
\newcommand{\rank}{\operatorname{rank}}
\newcommand{\im}{\operatorname{im}}
\newcommand{\sourcecommit}{\texttt{4eaee87548250b2682affda7d05d947c05554bae}}
\title{From Absolute Isolated-Zero Finiteness\\
to Uniform Bounds on Regular Fibers}
\author{Hilbert 16 Proof Audit project\\
\small AI-assisted mathematical exposition; independent review invited}
\date{2 November 2026}
\hypersetup{
 pdftitle={From Absolute Isolated-Zero Finiteness to Uniform Bounds on Regular Fibers},
 pdfauthor={Hilbert 16 Proof Audit project}}

\begin{document}
\maketitle

\begin{abstract}
We give a detailed conditional proof of the projection-counting principle
used in Section 11 of OpenAI's manuscript \emph{Uniform bounds for planar
polynomial limit cycles}. The hypotheses require a boundary-complete
finite list of domain margins and finiteness of the isolated solutions
of every fixed system in a specified differential and algebraic closure.
Under these hypotheses, a regular parameterized equation system has a
uniform bound on the number of connected components of its fibers.
The proof uses a proper function, generic linear perturbations, and
induction on the parameter dimension. A single perturbation is chosen
for a countable collection of witness fibers, avoiding both Fubini's
theorem and a separate argument at exceptional parameters. We spell out
the closure inheritance and the Jacobian calculation. Two counterexamples
show why neither finiteness for the original total zero set nor
finiteness for the full closure without boundary completeness suffices.
This note audits one conditional implication; it does not establish the
analytic hypotheses for polynomial-flow matching systems or certify
the full Hilbert-sixteenth manuscript.
\end{abstract}

{\small\tableofcontents}
\clearpage

\section{Scope, provenance, and the logical target}
\label{sec:scope}

The target is the implication
\[
 \begin{gathered}
 \text{absolute isolated-zero finiteness for a full finite closure}\\
 \Downarrow\\
 \text{uniform component bounds for regular fibers}.
 \end{gathered}
\]
It is Theorem 11.2 in the source manuscript \cite{OpenAI}, with its
hypotheses in Definition 11.1. We work from the public repository at the
fixed revision
\[
 \sourcecommit.
\]
The relevant source is
\href{https://github.com/openai/math/blob/4eaee87548250b2682affda7d05d947c05554bae/preprints/uniform-bounds-for-planar-polynomial-limit-cycles-September-24-2026/build/sections/11-counting.tex}{\texttt{build/sections/11-counting.tex}},
lines 22--217 at this revision. Pinning a revision records what is being
discussed; it is not a statement that other parts of that revision have
been checked.

The source uses a generic tilt for almost every parameter slice and then
extends the bound to exceptional parameters. The proof below has a
different organization. If the desired bound fails, choose a
\emph{countable} sequence of witness fibers. Sard's theorem allows one
tilt that is good on every member of this sequence. This is enough for
the contradiction and removes the exceptional-parameter step.

All mathematical arguments in this note are supplied for examination.
This document is AI-assisted, has not received a recorded independent
human referee report, and is not a formal proof-assistant certificate.
It makes no claim that the full Hilbert-sixteenth proof has been
verified. In particular, the source's Theorem 11.4, which asserts that
the required matching systems satisfy the closure hypothesis, is
\emph{not} an input proved here.

We use the following standard background facts:
\begin{enumerate}[label=(\roman*),leftmargin=*]
\item the implicit function theorem and the regular-level-set theorem;
\item Sard's theorem for smooth maps between finite-dimensional
second-countable manifolds;
\item the classification of connected one-dimensional manifolds without
boundary as circles or open intervals, up to diffeomorphism.
\end{enumerate}
These are background tools, not additional results about polynomial
flows. The final counterexample also uses the standard fact that a real
algebraic set has finitely many connected components. Sard's original
paper is listed in the references; conventional background sources
are listed for the manifold and semialgebraic facts
\cite{Sard,Lee,BCR}. No particular textbook theorem number is used.

\section{The hypotheses, with domains retained}
\label{sec:hypotheses}

Write the coordinates as
\[
 (p,x)\in\R^a\times\R^n.
\]
The letters \(p\) and \(x\) indicate their roles in a projection. They
do not indicate different kinds of analytic variables. In an absolute
zero test, the remaining parameter coordinates are also unknowns.

\begin{definition}[Open data and boundary-complete margins]
\label{def:data}
An open analytic datum consists of an open set
\(U\subset\R^a\times\R^n\), a real analytic map
\(F:U\to\R^m\), and a finite list of real analytic functions
\[
 g_1,\ldots,g_s:U\longrightarrow(0,\infty).
\]
The list is \emph{boundary-complete}, or satisfies
\emph{positive-margin completeness}, if the following holds:
whenever \(z_\nu\in U\) converges to a finite point \(z_*\) and
\[
 \inf_\nu g_j(z_\nu)>0\qquad(1\le j\le s),
\]
then \(z_*\in U\).
The equations are \(F=0\); the margins record its original open
restrictions, including coordinate-chart boundaries.
\end{definition}

Boundary completeness means precisely this sequential condition.
No compactness of \(U\), nor extension of \(F\) across a vanishing
margin, is assumed. A restriction \(h\ne0\) is recorded by the
positive margin \(h^2\). Empty zero sets and empty domains are allowed.

\begin{definition}[Finite generated closure]
\label{def:closure}
Starting from the coordinate functions, \(F_i\), \(g_j\), and fixed
real constants, permit the following finite operations:
\begin{enumerate}[label=(\alph*),leftmargin=*]
\item introduce finitely many additional independent real coordinates,
pulling back the old functions to the resulting product domain;
\item form finite sums and products, and take finitely many partial
derivatives in independent coordinates;
\item invert a generated function \(h\) on a domain where \(h\ne0\),
retaining that restriction as a margin;
\item append finitely many equations and positive margins made from
generated functions, while retaining every earlier equation and open
restriction.
\end{enumerate}
Coordinates may be regrouped as parameters and unknowns. Fixed real
values may be assigned to coordinates; equivalently, one retains those
coordinates and appends equations \(z_i=c_i\). Clearing a denominator
is permitted only while retaining the corresponding nonzero condition.

A \emph{descendant system} is any system obtained by finitely many such
operations, possibly repeated. Adjoining an arbitrary new analytic
function is not an allowed operation. Any analytic graph ties needed
in an application must already be among the data or be obtained by
the stated operations.
\end{definition}

Every particular descendant has finitely many coordinates, equations,
and margins. There is no common upper bound on their numbers or on the
derivative orders as the descendant varies. Fixed real constants are
unrestricted: the generic tilts selected below are allowed constants.

To avoid ambiguity about specialization, a partial derivative is taken
in an ambient independent coordinate before imposing an equation
fixing that coordinate. Retaining the fixing equation gives the same
solution set as eliminating the fixed coordinate, and preserves
isolation.

\begin{definition}[Absolute isolated-zero property for the closure]
\label{def:az}
The datum has property \(\AZ\) if, for every fixed descendant system,
its solution set has finitely many isolated points in its specified
open domain when \emph{all remaining coordinates} are treated as
unknowns. An isolated solution is isolated in the total solution set,
including all auxiliary coordinates and all retained equations.
\end{definition}

The quantifiers in this property are
\[
 \forall\text{ fixed finite descendants }E,\qquad
 \#\{\text{isolated total solutions of }E\}<\infty.
\]
There is no number required to work simultaneously for all \(E\).
The theorem will infer a uniform bound over \(p\) for \emph{one fixed
original system}. These are different assertions.

\begin{lemma}[Inheritance of the hypotheses]
\label{lem:inheritance}
Boundary completeness is preserved by introducing independent
coordinates, by a finite succession of generated strict restrictions
with their margins retained, and by taking fixed-coordinate slices.
Property \(\AZ\) is inherited by every descendant, with respect to
its own finite generated closure.
\end{lemma}

\begin{proof}
For additional coordinates, a convergent sequence projects to a
convergent sequence in the old domain. The original margins put the
limit in that domain, and the new coordinates have their given finite
limits. For a new condition \(h>0\), a positive lower bound on its
margin implies \(h(z_*)>0\) by continuity after the old domain has been
recovered. For \(h\ne0\), use the margin \(h^2\). Each new function is
analytic on the domain already constructed, so this argument applies
one stage at a time, including inverses.

For a slice obtained by fixing coordinates, append their fixed values
to the convergent sequence of the remaining coordinates. Apply the old
completeness condition and then take the slice again.

Finally, a finite sequence of permitted operations on a descendant is
still a finite sequence of permitted operations on the original datum.
All original equations and restrictions remain present. Thus the
absolute zero assertion for the original closure includes every system
in the descendant's closure.
\end{proof}

\section{The conditional uniformity theorem}
\label{sec:theorem}

\begin{theorem}[Uniform regular-fiber bound]
\label{thm:uniform}
Let \((U,F,g)\) satisfy Definitions \ref{def:data}--\ref{def:az}, with
\(0\le m\le n\). Suppose that
\[
 \rank F_x(p,x)=m
 \quad\text{at every }(p,x)\in U\text{ with }F(p,x)=0.
\]
Then there is a finite nonnegative integer \(N\), depending on this
fixed datum, such that every fiber
\[
 M_p=\{x:(p,x)\in U,\ F(p,x)=0\}
\]
has at most \(N\) connected components. If \(m=n\), every fiber has
at most \(N\) points.
\end{theorem}

For \(m=n\), the regular-level-set theorem makes each fiber discrete.
For \(m<n\), the fibers are smooth manifolds of dimension \(n-m\);
they need not be compact or connected. The theorem gives no effective
formula for \(N\). It makes no assertion about singular fibers in
which the displayed rank condition fails.

One can apply the theorem separately on the finitely many charts
specified by a nonzero \(m\times m\) minor of \(F_x\). The square of
the chosen minor is an additional margin. A connected component of
one chart's fiber belongs to only one component of the original
fiber, so summing the bounds for a finite covering bounds the original
components. Such a reduction is useful in applications; it does not
permit an infinite collection of charts.

The proof needs two mechanisms. A proper function gives a critical
point on every component. A common generic linear tilt makes those
critical points regular solutions of a finite multiplier system.
In the square case, allowing one parameter to vary first reduces the
number of parameters in that multiplier system.

\section{Proper functions and regular critical-point systems}
\label{sec:tools}

\subsection{A proper function from the margins}

\begin{lemma}[Proper barrier]
\label{lem:proper}
For boundary-complete open data define
\begin{equation}
 \rho(p,x)=1+|p|^2+|x|^2+\sum_{j=1}^s g_j(p,x)^{-2}.
 \label{eq:rho}
\end{equation}
On the equation locus, and on any fixed-coordinate slice of that
locus, \(\rho\) has compact sublevel sets. If \(z\) denotes the varying
coordinates of a slice, then for every fixed vector \(\ell\) the
function \(\rho+\ell\cdot z\) also has compact sublevel sets and is
bounded below. Its restriction to every nonempty connected component
of that slice attains a minimum.
\end{lemma}

\begin{proof}
A sublevel \(\rho\le C\) bounds all coordinates and gives
\(g_j\ge C^{-1/2}\) whenever the sublevel is nonempty. A sequence in
that sublevel has a convergent subsequence. Boundary completeness puts
its limit back in \(U\), and continuity preserves \(F=0\), the fixed
coordinates, and the sublevel inequality. The sublevel is therefore
compact in the ambient Euclidean space.

For a linear tilt, use
\[
 |z|^2+\ell\cdot z
 \ge \frac12|z|^2-\frac12|\ell|^2.
\]
All inverse-margin terms remain nonnegative. A tilted sublevel thus
still bounds the varying coordinates and keeps every margin bounded
away from zero. The fixed coordinates are constants, so the same
compactness argument applies. The inequality also gives a lower bound.

Let \(C_0\) be a nonempty connected component of a slice and choose
\(z_0\in C_0\). Components are closed in that slice. Intersecting \(C_0\)
with the compact sublevel at the value of \(z_0\) gives a nonempty
compact set on which the function attains a minimum. This minimum is
also the minimum on \(C_0\).
\end{proof}

All functions used in \eqref{eq:rho} belong to the generated closure.
At a later induction stage, the barrier is rebuilt using the squares
of \emph{all} coordinates then present, including multipliers, and the
inverse squares of all retained margins. Reusing a barrier that omitted
new unbounded coordinates would not establish the needed properness.

\subsection{The bordered Jacobian}

Let \(W\subset\R^q\) be open, let \(G:W\to\R^r\) be smooth, and
suppose \(DG\) has row rank \(r\) on \(M=\{G=0\}\). Let \(h:W\to\R\)
be smooth. At a constrained critical point of \(h+\ell\cdot z\), there
is a unique \(\lambda\in\R^r\) such that
\[
 \nabla h(z)+\ell=DG(z)^T\lambda.
\]
Existence follows because the gradient annihilates the tangent space
\(\ker DG(z)\); uniqueness follows from the injectivity of \(DG(z)^T\).

Define
\begin{equation}
 K_\ell(z,\lambda)=
 \begin{pmatrix}
 G(z)\\
 \nabla h(z)+\ell-DG(z)^T\lambda
 \end{pmatrix}.
 \label{eq:critical-general}
\end{equation}
There are \(r+q\) equations in \(q+r\) unknowns. At a zero set
\[
 A=DG(z),\qquad
 H=D_z^2\!\left(h+\ell\cdot z-\sum_{i=1}^r\lambda_iG_i\right).
\]
Direct differentiation gives the square matrix
\begin{equation}
 D_{(z,\lambda)}K_\ell=
 \begin{pmatrix}
 A&0\\
 H&-A^T
 \end{pmatrix}.
 \label{eq:border}
\end{equation}

\begin{lemma}[Regular multipliers and Morse critical points]
\label{lem:border}
The matrix \eqref{eq:border} is invertible if and only if the Hessian
of \((h+\ell\cdot z)|_M\) at the critical point is nondegenerate.
\end{lemma}

\begin{proof}
The constrained Hessian is the restriction of the symmetric bilinear
form \(H\) to \(T_zM=\ker A\). This follows by differentiating the
constraint along a local parametrization of \(M\): the second
derivatives of the parametrization in the Hessian of \(h+\ell\cdot z\) are
cancelled by the corresponding terms in
\(\sum_i\lambda_iD^2G_i\).

The kernel equations for \eqref{eq:border} are
\[
 Av=0,\qquad Hv=A^T\eta.
\]
If the restricted form is nondegenerate, testing the second equation
against every \(w\in\ker A\) gives \(v=0\). Then \(A^T\eta=0\), so
\(\eta=0\) because \(A\) has full row rank.

Conversely, if the restricted form is degenerate, there is a nonzero
\(v\in\ker A\) with \(w^THv=0\) for every \(w\in\ker A\). Hence
\[
 Hv\in(\ker A)^\perp=\im A^T.
\]
For some \(\eta\), \((v,\eta)\) is a nonzero kernel vector of
\eqref{eq:border}. This proves the equivalence.
\end{proof}

\subsection{A common generic tilt for countably many manifolds}

Here a function is \emph{Morse} when all its critical points are
nondegenerate. Distinct critical points need not have distinct values.

\begin{lemma}[Generic linear tilt]
\label{lem:tilt}
In the preceding setting, for almost every
\(\ell\in\R^q\), the function \((h+\ell\cdot z)|_M\) is Morse.
Consequently, for a countable collection of such pairs \((M_k,h_k)\)
in the same ambient dimension \(q\), a single \(\ell\) can be chosen
so that every \((h_k+\ell\cdot z)|_{M_k}\) is Morse.
No bound on the number of components of \(M_k\) is required.
\end{lemma}

\begin{proof}
The manifold \(M\times\R^r\) has dimension
\((q-r)+r=q\). Consider the smooth map
\[
 L:M\times\R^r\longrightarrow\R^q,\qquad
 L(z,\lambda)=DG(z)^T\lambda-\nabla h(z).
\]
Its fiber over \(\ell\) is precisely the critical-point system
\eqref{eq:critical-general}. At such a point, for a tangent vector
\((v,\eta)\) with \(Av=0\), its derivative is
\[
 DL(v,\eta)=-Hv+A^T\eta.
\]
Since source and target both have dimension \(q\), this derivative is
surjective exactly when its kernel is zero. By Lemma~\ref{lem:border},
that is exactly the nondegeneracy of the corresponding constrained
critical point.

Sard's theorem says that the critical values of \(L\) have Lebesgue
measure zero. Thus every regular value \(\ell\) makes all constrained
critical points nondegenerate. A value not in the image has no critical
points and satisfies the assertion vacuously.

The manifold is second countable: \(M\) is an embedded submanifold of
an open Euclidean set. Sard can therefore be applied through a
countable atlas even when there are infinitely many components.
For a countable collection, take the complement of the union of its
measure-zero bad sets. This complement has full measure, and any
single \(\ell\) in it has all the required properties.
\end{proof}

The lemma only concerns nondegeneracy. The existence of a minimum on
every component comes separately from Lemma~\ref{lem:proper}. Keeping
these two roles distinct is essential on noncompact manifolds.

\section{Square systems: induction on parameter dimension}
\label{sec:square}

\begin{lemma}[Monotonic parameter on a regular solution curve]
\label{lem:curve}
Suppose \(F(b,t,x)=0\) has \(n\) equations in \(x\in\R^n\), and
\(F_x\) is invertible on its zero set. For a fixed \(b\), let
\[
 C_b=\{(t,x):(b,t,x)\in U,\ F(b,t,x)=0\}.
\]
Every connected component of \(C_b\) meets a slice \(t=t_0\) in at
most one point.
\end{lemma}

\begin{proof}
The full row rank of \(F_{(t,x)}\) makes \(C_b\) a one-dimensional
embedded manifold without boundary. On its tangent space,
\[
 F_t\,dt+F_x\,dx=0.
\]
If \(dt=0\), then \(dx=0\); thus \(dt\) is nonzero on every nonzero
tangent vector.

A connected component is diffeomorphic to a circle or an open
interval. The circle case is impossible, since the real-valued
function \(t\) would attain a maximum and a minimum. In an interval
coordinate, the derivative of \(t\) is continuous and never zero,
so it has constant sign. Hence \(t\) is strictly monotone on the
component, which proves the claim.
\end{proof}

\begin{proposition}[Uniform point bound for square systems]
\label{prop:square}
Theorem~\ref{thm:uniform} holds when \(m=n\).
\end{proposition}

\begin{proof}
We induct on the number \(a\) of parameters. The induction statement
is simultaneous for all unknown dimensions \(n\) and for all data
satisfying the hypotheses. Increasing the number of unknowns in an
auxiliary system is therefore harmless; the quantity that decreases
is \(a\).

If \(n=0\), every fiber is empty or the unique point of \(\R^0\), so
the assertion is immediate. We may assume \(n\ge1\).

For \(a=0\), invertibility of \(F_x\) makes every solution isolated.
Property \(\AZ\), applied to the original system itself, gives the
required finiteness.

Assume the assertion proved for \(a-1\) parameters and every unknown
dimension. Write \(p=(b,t)\) with \(b\in\R^{a-1}\). If no uniform
bound existed, choose a sequence
\begin{equation}
 p_k=(b_k,t_k),\qquad \#M_{p_k}>k,\qquad k=1,2,\ldots.
 \label{eq:witness-square}
\end{equation}
There is no requirement that the parameters or their \(b\)-coordinates
be distinct. An infinite fiber can be used repeatedly. Equivalently,
one may select \(k+1\) distinct points from each chosen fiber.

For each \(k\), form the entire curve \(C_{b_k}\) of
Lemma~\ref{lem:curve}, allowing \(t\) to vary over its original
domain. That lemma implies that \(C_{b_k}\) has more than \(k\)
components: two of the selected points at the same \(t_k\) cannot
belong to the same curve component.

Apply Lemma~\ref{lem:tilt} to the countable collection
\[
 \left(C_{b_k},\ \rho(b_k,\cdot)\right),
 \qquad y=(t,x)\in\R^{n+1}.
\]
Choose one vector \(\ell\in\R^{n+1}\) that makes
\(\rho(b_k,y)+\ell\cdot y\) Morse on every \(C_{b_k}\).
This choice concerns each entire curve, even if it has infinitely
many components. By Lemma~\ref{lem:proper}, each component has a
minimum. Every such minimum is nondegenerate.

With this single \(\ell\) fixed, define
\begin{equation}
 K(b,y,\lambda)=
 \begin{pmatrix}
 F(b,y)\\
 \nabla_y\rho(b,y)+\ell-F_y(b,y)^T\lambda
 \end{pmatrix},\qquad \lambda\in\R^n.
 \label{eq:Ksquare}
\end{equation}
There are \(2n+1\) equations and \(2n+1\) unknowns
\(z=(y,\lambda)\), but only \(a-1\) parameters \(b\).
Let
\[
 \Delta(b,z)=\det D_zK(b,z),
\]
and use the open domain
\begin{equation}
 U_K=\{(b,t,x,\lambda):(b,t,x)\in U,\ \Delta(b,t,x,\lambda)\ne0\}.
 \label{eq:Kdomain}
\end{equation}
Retain the original margins and add \(\Delta^2\) as a margin.

We check the induction hypotheses rather than treating this as an
automatic reduction.
\begin{enumerate}[label=(\roman*),leftmargin=*]
\item \(K\) is a finite system in the generated closure.
The function \(\rho\) uses coordinates, sums, products, and inverses
of the positive margins. Its gradient uses partial derivatives;
the multiplier term uses additional coordinates and products.
The determinant \(\Delta\) is a polynomial in derivatives of these
expressions. The chosen \(\ell\) is a fixed real constant.
\item The new domain is boundary-complete with its retained margins.
This is Lemma~\ref{lem:inheritance}; in particular, a bounded sequence
whose new margin \(\Delta^2\) has a positive lower bound cannot leave
the regular Jacobian locus at its limit.
\item Property \(\AZ\) is inherited for the \emph{whole closure} of
this new system. Any later differentiation, inversion, multiplier
construction, or specialization is still a finite descendant of
the original system.
\item \(D_zK\) is invertible on \(U_K\), by its definition.
\end{enumerate}
The induction hypothesis therefore gives a finite bound \(N_K\)
on the number of solutions of \(K=0\) in \(U_K\), for every \(b\).

At a minimum on \(C_{b_k}\), the multiplier exists and is unique.
Lemma~\ref{lem:border} puts this solution in \(U_K\). Distinct
components give distinct \(y\)'s and hence distinct \((y,\lambda)\)'s.
Consequently the \(b_k\)-fiber of this \emph{one fixed} system has
more than \(k\) points. Taking \(k>N_K\) contradicts the induction
bound.
\end{proof}

\begin{remark}[No exceptional parameters are discarded]
\label{rem:exceptional}
The witness sequence in \eqref{eq:witness-square} is chosen before
the tilt. Its members may lie entirely in a lower-dimensional or
measure-zero parameter set. The common-tilt lemma treats those
particular slices directly. Thus the argument does not assume that
large fibers occur at generic parameters.
\end{remark}

\section{Positive-dimensional fibers}
\label{sec:components}

\begin{proof}[Proof of Theorem~\ref{thm:uniform}]
The square case is Proposition~\ref{prop:square}. Assume \(m<n\),
and suppose that no uniform component bound exists. Choose
parameters \(p_k\) such that \(M_{p_k}\) has more than \(k\)
connected components. Again repetitions are allowed.

Each \(M_{p_k}\) is a smooth manifold of dimension \(n-m\).
Lemma~\ref{lem:tilt}, with \(q=n\) and \(r=m\), provides a common
\(\ell\in\R^n\) such that
\[
 x\longmapsto\rho(p_k,x)+\ell\cdot x
\]
is Morse on every witness fiber. Its minimum on each component
exists by Lemma~\ref{lem:proper}.

For all parameters \(p\), use the same square multiplier system
\begin{equation}
 J(p,x,\lambda)=
 \begin{pmatrix}
 F(p,x)\\
 \nabla_x\rho(p,x)+\ell-F_x(p,x)^T\lambda
 \end{pmatrix}=0,\qquad \lambda\in\R^m.
 \label{eq:positive-dimensional}
\end{equation}
It has \(n+m\) equations in \(n+m\) unknowns \((x,\lambda)\).
Restrict to
\[
 \det D_{(x,\lambda)}J\ne0
\]
and retain the square of this determinant as a positive margin.
Exactly the checks in the proof of Proposition~\ref{prop:square}
show that this is admissible open data with inherited property
\(\AZ\). The square result, already proved for \emph{every}
parameter dimension \(a\), gives a uniform point bound for this
fixed system.

Every component of each witness fiber supplies a nondegenerate
minimum and therefore a solution in this regular locus.
Distinct components supply distinct solutions. There are more than
\(k\) such solutions at \(p_k\), a contradiction for large \(k\).
If \(m=0\), the construction has no multiplier coordinates; it
becomes \(\nabla_x\rho+\ell=0\), and the same proof applies.
\end{proof}

The two dimension reductions have different purposes. In the square
case a parameter is turned into a curve variable, reducing the
parameter count by one. For a positive-dimensional fiber, the
multiplier equations first make the problem square at the
\emph{same} parameter count. There is no circular invocation of the
general component theorem.

\section{Quantifiers and the scope of the conclusion}
\label{sec:quantifiers}

The following order is sufficient for the square-case contradiction:
\[
 \begin{gathered}
 \text{assume unbounded fiber cardinalities}\\
 \Longrightarrow\quad
 \text{choose all witnesses }(b_k,t_k)\\
 \Longrightarrow\
 \text{choose one }\ell\text{ good for every witness curve}\\
 \Longrightarrow\
 \text{form one finite system }K_\ell\\
 \Longrightarrow\
 \text{obtain one finite }N_{K_\ell}\\
 \Longrightarrow\
 \text{choose }k>N_{K_\ell}.
 \end{gathered}
\]
Although the chosen tilt can depend on the whole witness sequence,
the system contains only the finitely many entries of that vector
as constants. It does not contain the sequence as an infinite list
of equations. The bound is permitted to depend on the tilt; it is
still one fixed bound, contradicted by the infinitely many witnesses.

\begin{center}
\begin{tabular}{@{}p{0.57\linewidth}p{0.32\linewidth}@{}}
\toprule
Assertion & Role\\
\midrule
Each fixed finite descendant has finitely many isolated total solutions
& Required hypothesis\\[3pt]
A single bound for all descendants
& Not required\\[3pt]
One common tilt for countably many witness manifolds
& Supplied by Sard and countable intersection\\[3pt]
One common tilt good for every parameter value
& Not required\\[3pt]
A finite number of components before the tilt is chosen
& Not required\\[3pt]
All original open restrictions survive each construction
& Required throughout\\[3pt]
Uniform estimates over all derivative orders
& Not required\\[3pt]
An effective or computable component bound
& Not supplied\\
\bottomrule
\end{tabular}
\end{center}

At every induction step the new expressions involve only finitely
many derivatives. Parameter dimension decreases in the square
induction, so only finitely many such steps occur. This finiteness
does not weaken the universal closure assumption: one must know that
\(\AZ\) applies to the particular, possibly complicated descendants
created by those steps.

\begin{corollary}[The conditional interface to a cycle-counting argument]
\label{cor:interface}
Suppose a separate argument supplies finitely many fixed admissible
matching systems with property \(\AZ\), and represents every relevant
orbit, apart from at most \(T\) exceptions, by a point in one of
finitely many regular minor charts. Suppose also that distinct
represented orbits cannot belong to the same connected component
of a fiber in any one chart, at a parameter value common to the
orbits being counted. Then the number of those orbits is bounded
by \(T\) plus the sum of the component bounds of the charts.
\end{corollary}

\begin{proof}
Apply Theorem~\ref{thm:uniform} separately to the finitely many
charts. In one chart, different orbits occupy different components,
so its component bound bounds the number of orbits represented
there. Taking the union over charts can only overcount orbits.
Finally add the \(T\) exceptions.
\end{proof}

This corollary exposes the remaining obligation in the source
application. The existence of the finite matching description,
regularity, separation of orbit representations into components,
and property \(\AZ\) for \emph{every required finite descendant}
are additional mathematical assertions. They are not consequences
of Theorem~\ref{thm:uniform}. Nor does this note address the passage
from hyperbolic cycles to arbitrary isolated cycles or the handling
of the entire plane.

\section{Two counterexamples delimiting the hypotheses}
\label{sec:counterexamples}

\subsection{A: no isolated total zeros, but unbounded finite fibers}
\label{subsec:example-a}

Let
\[
 U=\{(p,x):0<p<1,\ 0<x<1\},\qquad
 F(p,x)=\sin(x/p).
\]
Take \(p\) as the parameter and \(x\) as the unknown. Use the complete
margin list
\[
 p,\quad1-p,\quad x,\quad1-x.
\]
If these margins stay positively bounded below along a convergent
sequence, its limit is still in \(U\). Thus the domain condition
holds. The function is real analytic there, and at every zero,
\[
 x=k\pi p,\qquad k\ge1,\qquad
 F_x(p,x)=\frac{(-1)^k}{p}\ne0.
\]
The fiber cardinality is
\[
 \#M_p=\#\{k\ge1:k\pi p<1\}
       =\left\lceil\frac1{\pi p}\right\rceil-1,
\]
which is finite for every \(p>0\) and unbounded as \(p\downarrow0\).

On the other hand, the total zero set has \emph{no isolated points}.
Each of its points lies on one of the open analytic curve segments
\[
 \{(p,k\pi p):0<p<1/(k\pi)\}.
\]
Finiteness of the isolated points of this total set is therefore
vacuously true.

The full closure detects the failure. Append the polynomial equation
\(x=1/2\). The resulting system has the infinitely many solutions
\[
 (p,x)=\left(\frac1{2\pi k},\frac12\right),
 \qquad k=1,2,\ldots.
\]
They are isolated. Indeed, the Jacobian determinant of
\((F,x-1/2)\) in \((p,x)\) is
\[
 F_p=-\frac{x}{p^2}\cos(x/p),
\]
which is nonzero at every displayed point. Thus property \(\AZ\)
fails for one very simple descendant.

This example preserves both the regularity and the boundary-complete
domain hypotheses. It disproves replacing Definition~\ref{def:az}
by finiteness of isolated points in the original total zero set,
even if every original parameter fiber is separately known to be
finite.

\subsection{B: isolated-zero finiteness without boundary completeness}
\label{subsec:example-b}

Let
\[
 D=\bigcup_{k\ge1}(3^{-k},\,2\cdot3^{-k}),\qquad
 U=D\times\R,\qquad F(x,y)=y.
\]
There are no parameters: \(a=0\), \(n=2\), and \(m=1\).
Use only the constant margin \(g_1=1\).
The derivative \(DF=(0,1)\) has row rank one everywhere. Its zero
set is
\[
 M=D\times\{0\},
\]
which has infinitely many connected components.

Nevertheless, this datum has property \(\AZ\) for the precise
generated closure of Definition~\ref{def:closure}. To see this,
start with the coordinates, \(F=y\), and \(g_1=1\). Every expression
generated by finitely many sums, products, partial derivatives,
additional coordinates, fixed real constants, and permitted
inverses is a rational function of the full coordinate tuple.
This follows inductively; in particular, a partial derivative of
a rational function is rational.

Take any fixed finite descendant. Clear its finitely many
denominators while retaining their nonzero restrictions. Its
equations become finitely many polynomial equations. Its permitted
domain is an open subset of the full Euclidean coordinate space:
it retains the original restriction \(x\in D\) and any finitely
many additional strict restrictions. Fixing a coordinate can be
represented by an additional polynomial equation.

If a solution is isolated after restricting that algebraic zero
set to this open domain, it is already an isolated point of the
algebraic zero set itself. In fact, a sufficiently small ball
about the point is contained in the domain, so any arbitrarily
nearby algebraic solutions would also be solutions of the
restricted system. A real algebraic set has finitely many
connected components, hence finitely many isolated points.
It follows that every such descendant has finitely many isolated
solutions. This verifies \(\AZ\), despite the infinitely many
components of \(M\).

The missing hypothesis is boundary completeness. For any fixed
\(k\), approach \(3^{-k}\) from inside
\((3^{-k},2\cdot3^{-k})\), keeping \(y=0\). The sequence is bounded
and the margin is always one, but its limit is outside \(U\).
The proposed barrier is
\[
 \rho(x,y)=2+x^2+y^2.
\]
On each component of \(M\), it has infimum \(2+3^{-2k}\) but no
minimum. Lemma~\ref{lem:proper} therefore fails exactly where the
counting argument needs it.

The restriction to the \emph{generated} closure matters in this
example. We have not claimed finiteness after adjoining arbitrary
new analytic functions on \(U\). Such an operation is absent from
Definition~\ref{def:closure}. The example shows that even the full
stated algebraic and differential closure cannot compensate for
an arbitrary open domain whose boundary is not recorded by its
margins.

\section{Review checkpoints and source comparison}
\label{sec:review}

A referee can check the conditional implication in the following
order:
\begin{enumerate}[leftmargin=*]
\item Verify that the domain condition makes the inverse-margin
barrier proper; inspect Lemma~\ref{lem:proper} and Example B.
\item Verify the square critical-point Jacobian
\eqref{eq:border}, including the uniqueness of multipliers.
\item Apply Sard to the map \(L\) in Lemma~\ref{lem:tilt}; check
that its source and target both have dimension \(q\).
\item Verify that \(t\) is globally injective on every solution-curve
component in Lemma~\ref{lem:curve}; local graph existence alone
would not justify that conclusion.
\item Follow the witness sequence through \eqref{eq:Ksquare}.
The critical system must be a single fixed system with \(a-1\)
parameters, not one separately chosen system for each \(k\).
\item Check all inherited margins and equations before using the
induction hypothesis.
\item Reduce general fibers to the square system
\eqref{eq:positive-dimensional}, at the unchanged parameter
dimension.
\end{enumerate}

\begin{center}
\begin{tabular}{@{}p{0.40\linewidth}p{0.49\linewidth}@{}}
\toprule
Pinned source, Section 11 & Exposition in this note\\
\midrule
Definition 11.1; lines 22--56
& Definitions \ref{def:data}--\ref{def:az}\\[3pt]
Proper function; lines 88--105
& Lemma \ref{lem:proper}\\[3pt]
Curve injectivity; lines 112--124
& Lemma \ref{lem:curve}\\[3pt]
Generic critical system; lines 126--159
& Lemmas \ref{lem:border} and \ref{lem:tilt}\\[3pt]
Parameter induction; lines 161--185
& Proposition \ref{prop:square}, using countable witnesses\\[3pt]
General fibers; lines 187--203
& Section \ref{sec:components}\\[3pt]
Closure warning; lines 207--216
& Section \ref{sec:quantifiers} and Example A\\
\bottomrule
\end{tabular}
\end{center}

The original proof's almost-everywhere parameter argument and its
implicit-function extension to exceptional parameters are not used
here. This alternate organization does not relax the substantive
closure hypothesis. A future audit of the full manuscript must
establish that hypothesis for the actual flow-passage systems and
their generated descendants, with their true open domains.

\begin{thebibliography}{9}

\bibitem{OpenAI}
OpenAI.
\emph{Uniform bounds for planar polynomial limit cycles}.
OpenAI Math Release preprint, manuscript dated 24 September 2026.
Section 11, Definition 11.1 and Theorem 11.2.
Repository revision \sourcecommit.
\href{https://github.com/openai/math/tree/4eaee87548250b2682affda7d05d947c05554bae/preprints/uniform-bounds-for-planar-polynomial-limit-cycles-September-24-2026}{Pinned manuscript directory}.

\bibitem{Sard}
Arthur Sard.
The measure of the critical values of differentiable maps.
\emph{Bulletin of the American Mathematical Society}
48 (1942), 883--890.
\href{https://doi.org/10.1090/S0002-9904-1942-07811-6}{doi:10.1090/S0002-9904-1942-07811-6}.

\bibitem{Lee}
John M. Lee.
\emph{Introduction to Smooth Manifolds}.
Second edition, Springer, 2013.
Background for regular level sets, Sard's theorem, and smooth
manifolds; no specific theorem number is used here.

\bibitem{BCR}
Jacek Bochnak, Michel Coste, and Marie-Fran\c{c}oise Roy.
\emph{Real Algebraic Geometry}.
Springer, 1998.
Background for finiteness of the connected components of a
semialgebraic set, used only in Example B.

\end{thebibliography}
\end{document}
